\documentclass[dvipdfmx, border=5pt]{standalone}

% フォント・数式関連
\usepackage[T1]{fontenc}
\usepackage{lmodern}
\usepackage{amsmath, amssymb}

% TikZ関連
\usepackage{tikz}
\usetikzlibrary{shapes.geometric, arrows.meta, positioning, calc}

% スタイル定義（色味は維持しつつシンプルに）
\tikzset{
    base/.style={rectangle, draw, rounded corners, inner sep=8pt, align=center, font=\sffamily\small, line width=0.8pt},
    user/.style={base, fill=gray!10, draw=gray!70!black},
    llm/.style={base, fill=blue!10, draw=blue!70!black, minimum width=3cm},
    mcp/.style={base, fill=purple!10, draw=purple!70!black, minimum width=3cm},
    blender/.style={base, fill=green!10, draw=green!70!black, minimum width=3cm},
    storage/.style={base, fill=orange!10, draw=orange!70!black},
    arrow/.style={-{Stealth[length=6pt, width=5pt]}, line width=0.8pt, draw=black!80},
    darrow/.style={{Stealth[length=6pt, width=5pt]}-{Stealth[length=6pt, width=5pt]}, line width=0.8pt, draw=black!80},
    dashed_arrow/.style={-{Stealth[length=6pt, width=5pt]}, dashed, line width=0.8pt, draw=black!80}
}

\begin{document}

\begin{tikzpicture}[node distance=1.2cm]

    % --- ノード配置 (上から下へ) ---
    
    % 1. ユーザ層
    \node (User) [user] {ユーザ / プロンプト};

    % 2. AIコア層 (LLM + Planner)
    \node (AICore) [llm, below=of User] {AIコア\\(LLM \& Planner)};

    % 3. MCP統合層 (Router + Tools)
    \node (MCP) [mcp, below=of AICore] {MCP統合システム\\(Router \& Tools)};

    % 4. Blender環境層 (API + Scene)
    \node (Blender) [blender, below=of MCP] {Blender 3D環境\\(Python API \& GUI)};

    % 5. 知識層 (DB) - サイドに配置
    \node (Knowledge) [storage, right=1.5cm of MCP] {知識ベース\\(Static/Dynamic DB)};

    % --- 接続 ---

    % ユーザ -> AI
    \draw [arrow] (User) -- (AICore);

    % AI <-> MCP (命令と結果)
    \draw [darrow] (AICore) -- (MCP);

    % MCP -> Blender (実行)
    \draw [darrow] (MCP) -- node[right, font=\scriptsize] {API操作} (Blender);

    % DBとの連携 (MCPがツール経由でアクセス、あるいはAIが参照)
    \draw [darrow] (MCP) -- (Knowledge);
    \draw [dashed_arrow] (Knowledge) |- (AICore);

\end{tikzpicture}

\end{document}